% --- Archon physics formalization source begin ---
% archon:physics
% archon:covers PhyXMiniProblems/problem_phyx_mini_0582.lean
% archon:source-report reports/phyx_mini/problem_phyx_mini_0582.source.json

\chapter{Physics problem phyx\_mini\_0582}
\label{ch:PhyXMiniProblems_problem_phyx_mini_0582}

\paragraph{Problem source.}
\#\# Physical scenario

A diatomic gas molecule consists of two atoms of mass \$m\$ separated by a fixed distance \$d\$ rotating about an axis as indicated in figure.

\#\# Question

Assuming that its angular momentum is quantized as in the Bohr model for the hydrogen atom, find the possible quantized rotational energies.

\#\# Answer choices

- A: A: \textbackslash{}frac\{n\textasciicircum{}2 h\textasciicircum{}2\}\{2 \textbackslash{}pi\textasciicircum{}2 m d\textasciicircum{}2\}

- B: B: \textbackslash{}frac\{n\textasciicircum{}2 h\textasciicircum{}2\}\{8 \textbackslash{}pi\textasciicircum{}2 m d\textasciicircum{}2\}

- C: C: \textbackslash{}frac\{n\textasciicircum{}2 h\textasciicircum{}2\}\{4 \textbackslash{}pi\textasciicircum{}2 m d\textasciicircum{}2\}

- D: D: \textbackslash{}frac\{n h\textasciicircum{}2\}\{4 \textbackslash{}pi\textasciicircum{}2 m d\textasciicircum{}2\}

\#\# Figure caption (auxiliary; use the image as primary evidence)

The image shows a diagram of a rotating system around an axis. Here are the details:

- **Axis**: A vertical line labeled "Axis" represents the axis of rotation.
- **Angular Velocity (\textbackslash{}(\textbackslash{}omega\textbackslash{}))**: An arrow curving around the axis indicates the direction of rotation, labeled with the Greek letter \textbackslash{}(\textbackslash{}omega\textbackslash{}).
- **Masses**: Two spherical objects, labeled with "m", are positioned on a circular path, implying they have equal mass.
- **Circular Path**: A dashed circle outlines the path of the masses, showing they rotate symmetrically around the axis.
- **Distances**: The distance from each mass to the axis is labeled as \textbackslash{}(d/2\textbackslash{}).
- **Text**: The diagram includes labels "Axis", "\textbackslash{}(\textbackslash{}omega\textbackslash{})", "m", and "d/2".

This system appears to illustrate a simple model in rotational dynamics, possibly representing a two-mass system rotating around a central axis.

\#\# Dataset metadata

Quantum Phenomena; Physical Model Grounding Reasoning, Numerical Reasoning

\paragraph{Recorded answer/context.}
C: C: \textbackslash{}frac\{n\textasciicircum{}2 h\textasciicircum{}2\}\{4 \textbackslash{}pi\textasciicircum{}2 m d\textasciicircum{}2\}

\paragraph{Figure/image path.}
/root/proposal\_for\_physic/science-mango/phyx\_mini\_run/phyx\_data/test\_image/582.png

\paragraph{Formalization target.}
create a compiling Lean file with sorry bodies at `PhyXMiniProblems/problem\_phyx\_mini\_0582.lean`. The Lean declarations must preserve the physical quantities, dimensions or dimensional roles, figure labels, governing-law hypotheses, and final relation expressed by this problem.
Use Mathlib/Physlib names found through LeanExplore where available. If a physics API is missing, introduce faithful local abstractions rather than scalar placeholder aliases.

\begin{theorem}[Physics formalization target]
\label{thm:physics:phyx_mini_0582:target}
\lean{PhyXMiniProblems.ProblemPhyXMini0582.problem_phyx_mini_0582}
\uses{def:physics:phyx-mini-0582:phyxminiproblems-problemphyxmini0582-massreadout, def:physics:phyx-mini-0582:phyxminiproblems-problemphyxmini0582-lengthreadout, def:physics:phyx-mini-0582:phyxminiproblems-problemphyxmini0582-planckconstantreadout, def:physics:phyx-mini-0582:phyxminiproblems-problemphyxmini0582-energyreadout, def:physics:phyx-mini-0582:phyxminiproblems-problemphyxmini0582-diatomicmolecularrotorsetup, def:physics:phyx-mini-0582:phyxminiproblems-problemphyxmini0582-matchesdiatomicrotorscenario, def:physics:phyx-mini-0582:phyxminiproblems-problemphyxmini0582-matchessupplieddiatomicrotorfigure, def:physics:phyx-mini-0582:phyxminiproblems-problemphyxmini0582-hasphysicaldiatomicrotorparameters, def:physics:phyx-mini-0582:phyxminiproblems-problemphyxmini0582-satisfiescentereddiatomicgeometry, def:physics:phyx-mini-0582:phyxminiproblems-problemphyxmini0582-satisfiesrigiddiatomicrotorlaws, def:physics:phyx-mini-0582:phyxminiproblems-problemphyxmini0582-satisfiesbohrangularmomentumquantization, def:physics:phyx-mini-0582:phyxminiproblems-problemphyxmini0582-recordeddatasetanswer, def:physics:phyx-mini-0582:phyxminiproblems-problemphyxmini0582-matchesanswerchoice}
The assigned autoformalize agent should translate this physics problem into Lean declarations in the covered file, with theorem and lemma proof bodies written as `by sorry`.
\end{theorem}
\begin{proof}
This is an autoformalization task, not a proof task. Produce statements that can later be proved by the physics prover without weakening the physical model.
\end{proof}

% --- Archon declaration topology begin ---
\section{Declaration topology}
\begin{definition}[angular Speed Dimension]
  \label{def:physics:phyx-mini-0582:phyxminiproblems-problemphyxmini0582-angularspeeddimension}
  \lean{PhyXMiniProblems.ProblemPhyXMini0582.angularSpeedDimension}
  Angular speed has physical dimension inverse time
\end{definition}
\begin{proof}
  This is the declared model definition of angular Speed Dimension.
\end{proof}
\begin{definition}[moment Of Inertia Dimension]
  \label{def:physics:phyx-mini-0582:phyxminiproblems-problemphyxmini0582-momentofinertiadimension}
  \lean{PhyXMiniProblems.ProblemPhyXMini0582.momentOfInertiaDimension}
  Moment of inertia has physical dimension mass times length squared
\end{definition}
\begin{proof}
  This is the declared model definition of moment Of Inertia Dimension.
\end{proof}
\begin{definition}[action Dimension]
  \label{def:physics:phyx-mini-0582:phyxminiproblems-problemphyxmini0582-actiondimension}
  \lean{PhyXMiniProblems.ProblemPhyXMini0582.actionDimension}
  Action and angular momentum have dimension mass times length squared per time
\end{definition}
\begin{proof}
  This is the declared model definition of action Dimension.
\end{proof}
\begin{definition}[Mass Quantity]
  \label{def:physics:phyx-mini-0582:phyxminiproblems-problemphyxmini0582-massquantity}
  \lean{PhyXMiniProblems.ProblemPhyXMini0582.MassQuantity}
  A nonnegative, unit-independent physical mass
\end{definition}
\begin{proof}
  This is the declared model definition of Mass Quantity.
\end{proof}
\begin{definition}[Length Quantity]
  \label{def:physics:phyx-mini-0582:phyxminiproblems-problemphyxmini0582-lengthquantity}
  \lean{PhyXMiniProblems.ProblemPhyXMini0582.LengthQuantity}
  A nonnegative, unit-independent physical length
\end{definition}
\begin{proof}
  This is the declared model definition of Length Quantity.
\end{proof}
\begin{definition}[Angular Speed Quantity]
  \label{def:physics:phyx-mini-0582:phyxminiproblems-problemphyxmini0582-angularspeedquantity}
  \lean{PhyXMiniProblems.ProblemPhyXMini0582.AngularSpeedQuantity}
  \uses{def:physics:phyx-mini-0582:phyxminiproblems-problemphyxmini0582-angularspeeddimension}
  A nonnegative angular-speed magnitude
\end{definition}
\begin{proof}
  This is the declared model definition of Angular Speed Quantity.
\end{proof}
\begin{definition}[Moment Of Inertia Quantity]
  \label{def:physics:phyx-mini-0582:phyxminiproblems-problemphyxmini0582-momentofinertiaquantity}
  \lean{PhyXMiniProblems.ProblemPhyXMini0582.MomentOfInertiaQuantity}
  \uses{def:physics:phyx-mini-0582:phyxminiproblems-problemphyxmini0582-momentofinertiadimension}
  A nonnegative moment of inertia about the depicted axis
\end{definition}
\begin{proof}
  This is the declared model definition of Moment Of Inertia Quantity.
\end{proof}
\begin{definition}[Angular Momentum Magnitude]
  \label{def:physics:phyx-mini-0582:phyxminiproblems-problemphyxmini0582-angularmomentummagnitude}
  \lean{PhyXMiniProblems.ProblemPhyXMini0582.AngularMomentumMagnitude}
  \uses{def:physics:phyx-mini-0582:phyxminiproblems-problemphyxmini0582-actiondimension}
  A nonnegative angular-momentum magnitude
\end{definition}
\begin{proof}
  This is the declared model definition of Angular Momentum Magnitude.
\end{proof}
\begin{definition}[Planck Constant Quantity]
  \label{def:physics:phyx-mini-0582:phyxminiproblems-problemphyxmini0582-planckconstantquantity}
  \lean{PhyXMiniProblems.ProblemPhyXMini0582.PlanckConstantQuantity}
  \uses{def:physics:phyx-mini-0582:phyxminiproblems-problemphyxmini0582-actiondimension}
  The ordinary Planck constant h, as a unit-independent action
\end{definition}
\begin{proof}
  This is the declared model definition of Planck Constant Quantity.
\end{proof}
\begin{definition}[mass Readout]
  \label{def:physics:phyx-mini-0582:phyxminiproblems-problemphyxmini0582-massreadout}
  \lean{PhyXMiniProblems.ProblemPhyXMini0582.massReadout}
  \uses{def:physics:phyx-mini-0582:phyxminiproblems-problemphyxmini0582-massquantity}
  Read a physical mass in a selected mass unit
\end{definition}
\begin{proof}
  This is the declared model definition of mass Readout.
\end{proof}
\begin{definition}[length Readout]
  \label{def:physics:phyx-mini-0582:phyxminiproblems-problemphyxmini0582-lengthreadout}
  \lean{PhyXMiniProblems.ProblemPhyXMini0582.lengthReadout}
  \uses{def:physics:phyx-mini-0582:phyxminiproblems-problemphyxmini0582-lengthquantity}
  Read a physical length in a selected length unit
\end{definition}
\begin{proof}
  This is the declared model definition of length Readout.
\end{proof}
\begin{definition}[angular Speed Readout]
  \label{def:physics:phyx-mini-0582:phyxminiproblems-problemphyxmini0582-angularspeedreadout}
  \lean{PhyXMiniProblems.ProblemPhyXMini0582.angularSpeedReadout}
  \uses{def:physics:phyx-mini-0582:phyxminiproblems-problemphyxmini0582-angularspeedquantity}
  Read angular speed in inverse units of the selected time unit
\end{definition}
\begin{proof}
  This is the declared model definition of angular Speed Readout.
\end{proof}
\begin{definition}[moment Of Inertia Readout]
  \label{def:physics:phyx-mini-0582:phyxminiproblems-problemphyxmini0582-momentofinertiareadout}
  \lean{PhyXMiniProblems.ProblemPhyXMini0582.momentOfInertiaReadout}
  \uses{def:physics:phyx-mini-0582:phyxminiproblems-problemphyxmini0582-momentofinertiaquantity}
  Read moment of inertia in coherent selected mass and length units
\end{definition}
\begin{proof}
  This is the declared model definition of moment Of Inertia Readout.
\end{proof}
\begin{definition}[angular Momentum Readout]
  \label{def:physics:phyx-mini-0582:phyxminiproblems-problemphyxmini0582-angularmomentumreadout}
  \lean{PhyXMiniProblems.ProblemPhyXMini0582.angularMomentumReadout}
  \uses{def:physics:phyx-mini-0582:phyxminiproblems-problemphyxmini0582-angularmomentummagnitude}
  Read angular momentum in coherent mechanical units
\end{definition}
\begin{proof}
  This is the declared model definition of angular Momentum Readout.
\end{proof}
\begin{definition}[planck Constant Readout]
  \label{def:physics:phyx-mini-0582:phyxminiproblems-problemphyxmini0582-planckconstantreadout}
  \lean{PhyXMiniProblems.ProblemPhyXMini0582.planckConstantReadout}
  \uses{def:physics:phyx-mini-0582:phyxminiproblems-problemphyxmini0582-planckconstantquantity}
  Read h in the same action units as angular momentum
\end{definition}
\begin{proof}
  This is the declared model definition of planck Constant Readout.
\end{proof}
\begin{definition}[energy Readout]
  \label{def:physics:phyx-mini-0582:phyxminiproblems-problemphyxmini0582-energyreadout}
  \lean{PhyXMiniProblems.ProblemPhyXMini0582.energyReadout}
  Read energy in the coherent unit induced by the selected base units
\end{definition}
\begin{proof}
  This is the declared model definition of energy Readout.
\end{proof}
\begin{definition}[Atom Role]
  \label{def:physics:phyx-mini-0582:phyxminiproblems-problemphyxmini0582-atomrole}
  \lean{PhyXMiniProblems.ProblemPhyXMini0582.AtomRole}
  The two equal atoms on opposite sides of the rotation axis
\end{definition}
\begin{proof}
  This is the declared model definition of Atom Role.
\end{proof}
\begin{definition}[Molecular Rotor Model]
  \label{def:physics:phyx-mini-0582:phyxminiproblems-problemphyxmini0582-molecularrotormodel}
  \lean{PhyXMiniProblems.ProblemPhyXMini0582.MolecularRotorModel}
  The physical approximation used for the molecule
\end{definition}
\begin{proof}
  This is the declared model definition of Molecular Rotor Model.
\end{proof}
\begin{definition}[Axis Geometry]
  \label{def:physics:phyx-mini-0582:phyxminiproblems-problemphyxmini0582-axisgeometry}
  \lean{PhyXMiniProblems.ProblemPhyXMini0582.AxisGeometry}
  The relation between the depicted axis and the internuclear bond
\end{definition}
\begin{proof}
  This is the declared model definition of Axis Geometry.
\end{proof}
\begin{definition}[Figure Object]
  \label{def:physics:phyx-mini-0582:phyxminiproblems-problemphyxmini0582-figureobject}
  \lean{PhyXMiniProblems.ProblemPhyXMini0582.FigureObject}
  Objects and graphical marks visible in the supplied raster
\end{definition}
\begin{proof}
  This is the declared model definition of Figure Object.
\end{proof}
\begin{definition}[Figure Text Label]
  \label{def:physics:phyx-mini-0582:phyxminiproblems-problemphyxmini0582-figuretextlabel}
  \lean{PhyXMiniProblems.ProblemPhyXMini0582.FigureTextLabel}
  Individual occurrences of text visible in the supplied raster
\end{definition}
\begin{proof}
  This is the declared model definition of Figure Text Label.
\end{proof}
\begin{definition}[Supplied Diatomic Rotor Figure]
  \label{def:physics:phyx-mini-0582:phyxminiproblems-problemphyxmini0582-supplieddiatomicrotorfigure}
  \lean{PhyXMiniProblems.ProblemPhyXMini0582.SuppliedDiatomicRotorFigure}
  \uses{def:physics:phyx-mini-0582:phyxminiproblems-problemphyxmini0582-atomrole, def:physics:phyx-mini-0582:phyxminiproblems-problemphyxmini0582-figureobject, def:physics:phyx-mini-0582:phyxminiproblems-problemphyxmini0582-figuretextlabel}
  Literal and qualitative information read from image 582. The drawn orbit is an ellipse only because a circular orbit is displayed in perspective
\end{definition}
\begin{proof}
  This is the declared model definition of Supplied Diatomic Rotor Figure.
\end{proof}
\begin{definition}[Diatomic Molecular Rotor Setup]
  \label{def:physics:phyx-mini-0582:phyxminiproblems-problemphyxmini0582-diatomicmolecularrotorsetup}
  \lean{PhyXMiniProblems.ProblemPhyXMini0582.DiatomicMolecularRotorSetup}
  \uses{def:physics:phyx-mini-0582:phyxminiproblems-problemphyxmini0582-massquantity, def:physics:phyx-mini-0582:phyxminiproblems-problemphyxmini0582-lengthquantity, def:physics:phyx-mini-0582:phyxminiproblems-problemphyxmini0582-angularspeedquantity, def:physics:phyx-mini-0582:phyxminiproblems-problemphyxmini0582-momentofinertiaquantity, def:physics:phyx-mini-0582:phyxminiproblems-problemphyxmini0582-angularmomentummagnitude, def:physics:phyx-mini-0582:phyxminiproblems-problemphyxmini0582-planckconstantquantity, def:physics:phyx-mini-0582:phyxminiproblems-problemphyxmini0582-atomrole, def:physics:phyx-mini-0582:phyxminiproblems-problemphyxmini0582-molecularrotormodel, def:physics:phyx-mini-0582:phyxminiproblems-problemphyxmini0582-axisgeometry, def:physics:phyx-mini-0582:phyxminiproblems-problemphyxmini0582-supplieddiatomicrotorfigure}
  Independent quantities and roles of the rotating molecule. In particular, the inertia and level energies are observables to be related by laws, not definitions containing the requested closed form
\end{definition}
\begin{proof}
  This is the declared model definition of Diatomic Molecular Rotor Setup.
\end{proof}
\begin{definition}[Matches Diatomic Rotor Scenario]
  \label{def:physics:phyx-mini-0582:phyxminiproblems-problemphyxmini0582-matchesdiatomicrotorscenario}
  \lean{PhyXMiniProblems.ProblemPhyXMini0582.MatchesDiatomicRotorScenario}
  \uses{def:physics:phyx-mini-0582:phyxminiproblems-problemphyxmini0582-massreadout, def:physics:phyx-mini-0582:phyxminiproblems-problemphyxmini0582-diatomicmolecularrotorsetup}
  The prose-level rigid, equal-mass diatomic-rotor scenario
\end{definition}
\begin{proof}
  This is the declared model definition of Matches Diatomic Rotor Scenario.
\end{proof}
\begin{definition}[Matches Supplied Diatomic Rotor Figure]
  \label{def:physics:phyx-mini-0582:phyxminiproblems-problemphyxmini0582-matchessupplieddiatomicrotorfigure}
  \lean{PhyXMiniProblems.ProblemPhyXMini0582.MatchesSuppliedDiatomicRotorFigure}
  \uses{def:physics:phyx-mini-0582:phyxminiproblems-problemphyxmini0582-diatomicmolecularrotorsetup}
  Exact labels and qualitative geometry transcribed from the primary image
\end{definition}
\begin{proof}
  This is the declared model definition of Matches Supplied Diatomic Rotor Figure.
\end{proof}
\begin{definition}[Has Physical Diatomic Rotor Parameters]
  \label{def:physics:phyx-mini-0582:phyxminiproblems-problemphyxmini0582-hasphysicaldiatomicrotorparameters}
  \lean{PhyXMiniProblems.ProblemPhyXMini0582.HasPhysicalDiatomicRotorParameters}
  \uses{def:physics:phyx-mini-0582:phyxminiproblems-problemphyxmini0582-massreadout, def:physics:phyx-mini-0582:phyxminiproblems-problemphyxmini0582-lengthreadout, def:physics:phyx-mini-0582:phyxminiproblems-problemphyxmini0582-momentofinertiareadout, def:physics:phyx-mini-0582:phyxminiproblems-problemphyxmini0582-planckconstantreadout, def:physics:phyx-mini-0582:phyxminiproblems-problemphyxmini0582-energyreadout, def:physics:phyx-mini-0582:phyxminiproblems-problemphyxmini0582-atomrole, def:physics:phyx-mini-0582:phyxminiproblems-problemphyxmini0582-diatomicmolecularrotorsetup}
  Positivity and nonnegativity conditions for the physical magnitudes
\end{definition}
\begin{proof}
  This is the declared model definition of Has Physical Diatomic Rotor Parameters.
\end{proof}
\begin{definition}[Satisfies Centered Diatomic Geometry]
  \label{def:physics:phyx-mini-0582:phyxminiproblems-problemphyxmini0582-satisfiescentereddiatomicgeometry}
  \lean{PhyXMiniProblems.ProblemPhyXMini0582.SatisfiesCenteredDiatomicGeometry}
  \uses{def:physics:phyx-mini-0582:phyxminiproblems-problemphyxmini0582-lengthreadout, def:physics:phyx-mini-0582:phyxminiproblems-problemphyxmini0582-atomrole, def:physics:phyx-mini-0582:phyxminiproblems-problemphyxmini0582-diatomicmolecularrotorsetup}
  Each atom is at the figure-labelled radius d/2 from the centered axis
\end{definition}
\begin{proof}
  This is the declared model definition of Satisfies Centered Diatomic Geometry.
\end{proof}
\begin{definition}[Satisfies Rigid Diatomic Rotor Laws]
  \label{def:physics:phyx-mini-0582:phyxminiproblems-problemphyxmini0582-satisfiesrigiddiatomicrotorlaws}
  \lean{PhyXMiniProblems.ProblemPhyXMini0582.SatisfiesRigidDiatomicRotorLaws}
  \uses{def:physics:phyx-mini-0582:phyxminiproblems-problemphyxmini0582-massreadout, def:physics:phyx-mini-0582:phyxminiproblems-problemphyxmini0582-lengthreadout, def:physics:phyx-mini-0582:phyxminiproblems-problemphyxmini0582-angularspeedreadout, def:physics:phyx-mini-0582:phyxminiproblems-problemphyxmini0582-momentofinertiareadout, def:physics:phyx-mini-0582:phyxminiproblems-problemphyxmini0582-angularmomentumreadout, def:physics:phyx-mini-0582:phyxminiproblems-problemphyxmini0582-energyreadout, def:physics:phyx-mini-0582:phyxminiproblems-problemphyxmini0582-atomrole, def:physics:phyx-mini-0582:phyxminiproblems-problemphyxmini0582-diatomicmolecularrotorsetup}
  Classical laws for a two-point-mass rotor: I = Σ mᵢrᵢ², L = Iω, and E = Iω²/2. None contains the requested expression in m, d, h, and n
\end{definition}
\begin{proof}
  This is the declared model definition of Satisfies Rigid Diatomic Rotor Laws.
\end{proof}
\begin{definition}[Satisfies Bohr Angular Momentum Quantization]
  \label{def:physics:phyx-mini-0582:phyxminiproblems-problemphyxmini0582-satisfiesbohrangularmomentumquantization}
  \lean{PhyXMiniProblems.ProblemPhyXMini0582.SatisfiesBohrAngularMomentumQuantization}
  \uses{def:physics:phyx-mini-0582:phyxminiproblems-problemphyxmini0582-angularmomentumreadout, def:physics:phyx-mini-0582:phyxminiproblems-problemphyxmini0582-planckconstantreadout, def:physics:phyx-mini-0582:phyxminiproblems-problemphyxmini0582-diatomicmolecularrotorsetup}
  Bohr quantization of angular-momentum magnitude: Lₙ = n h/(2π)
\end{definition}
\begin{proof}
  This is the declared model definition of Satisfies Bohr Angular Momentum Quantization.
\end{proof}
\begin{lemma}[diatomic moment Of Inertia readout]
  \label{lem:physics:phyx-mini-0582:phyxminiproblems-problemphyxmini0582-diatomic-momentofinertia-readout}
  \lean{PhyXMiniProblems.ProblemPhyXMini0582.diatomic_momentOfInertia_readout}
  \uses{def:physics:phyx-mini-0582:phyxminiproblems-problemphyxmini0582-massreadout, def:physics:phyx-mini-0582:phyxminiproblems-problemphyxmini0582-lengthreadout, def:physics:phyx-mini-0582:phyxminiproblems-problemphyxmini0582-momentofinertiareadout, def:physics:phyx-mini-0582:phyxminiproblems-problemphyxmini0582-diatomicmolecularrotorsetup, def:physics:phyx-mini-0582:phyxminiproblems-problemphyxmini0582-matchesdiatomicrotorscenario, def:physics:phyx-mini-0582:phyxminiproblems-problemphyxmini0582-satisfiescentereddiatomicgeometry, def:physics:phyx-mini-0582:phyxminiproblems-problemphyxmini0582-satisfiesrigiddiatomicrotorlaws}
  Two equal masses at radius d/2 have I = md²/2
\end{lemma}
\begin{proof}
  The conclusion follows from the governing relations and figure readouts named in the statement of diatomic moment Of Inertia readout.
\end{proof}
\begin{lemma}[rotational Energy eq angular Momentum Squared div inertia]
  \label{lem:physics:phyx-mini-0582:phyxminiproblems-problemphyxmini0582-rotationalenergy-eq-angularmomentumsquared-div-inertia}
  \lean{PhyXMiniProblems.ProblemPhyXMini0582.rotationalEnergy_eq_angularMomentumSquared_div_inertia}
  \uses{def:physics:phyx-mini-0582:phyxminiproblems-problemphyxmini0582-momentofinertiareadout, def:physics:phyx-mini-0582:phyxminiproblems-problemphyxmini0582-angularmomentumreadout, def:physics:phyx-mini-0582:phyxminiproblems-problemphyxmini0582-energyreadout, def:physics:phyx-mini-0582:phyxminiproblems-problemphyxmini0582-diatomicmolecularrotorsetup, def:physics:phyx-mini-0582:phyxminiproblems-problemphyxmini0582-hasphysicaldiatomicrotorparameters, def:physics:phyx-mini-0582:phyxminiproblems-problemphyxmini0582-satisfiesrigiddiatomicrotorlaws}
  The scalar-axis rotor laws imply E = L²/(2I)
\end{lemma}
\begin{proof}
  The conclusion follows from the governing relations and figure readouts named in the statement of rotational Energy eq angular Momentum Squared div inertia.
\end{proof}
\begin{lemma}[quantized rotational Energy readout]
  \label{lem:physics:phyx-mini-0582:phyxminiproblems-problemphyxmini0582-quantized-rotationalenergy-readout}
  \lean{PhyXMiniProblems.ProblemPhyXMini0582.quantized_rotationalEnergy_readout}
  \uses{def:physics:phyx-mini-0582:phyxminiproblems-problemphyxmini0582-massreadout, def:physics:phyx-mini-0582:phyxminiproblems-problemphyxmini0582-lengthreadout, def:physics:phyx-mini-0582:phyxminiproblems-problemphyxmini0582-planckconstantreadout, def:physics:phyx-mini-0582:phyxminiproblems-problemphyxmini0582-energyreadout, def:physics:phyx-mini-0582:phyxminiproblems-problemphyxmini0582-diatomicmolecularrotorsetup, def:physics:phyx-mini-0582:phyxminiproblems-problemphyxmini0582-matchesdiatomicrotorscenario, def:physics:phyx-mini-0582:phyxminiproblems-problemphyxmini0582-hasphysicaldiatomicrotorparameters, def:physics:phyx-mini-0582:phyxminiproblems-problemphyxmini0582-satisfiescentereddiatomicgeometry, def:physics:phyx-mini-0582:phyxminiproblems-problemphyxmini0582-satisfiesrigiddiatomicrotorlaws, def:physics:phyx-mini-0582:phyxminiproblems-problemphyxmini0582-satisfiesbohrangularmomentumquantization}
  The possible rotational energies are Eₙ = n² h² / (4 π² m d²) in every coherent unit system
\end{lemma}
\begin{proof}
  The conclusion follows from the governing relations and figure readouts named in the statement of quantized rotational Energy readout.
\end{proof}
\begin{definition}[Answer Choice]
  \label{def:physics:phyx-mini-0582:phyxminiproblems-problemphyxmini0582-answerchoice}
  \lean{PhyXMiniProblems.ProblemPhyXMini0582.AnswerChoice}
  Labels of the four formulas displayed in the source
\end{definition}
\begin{proof}
  This is the declared model definition of Answer Choice.
\end{proof}
\begin{definition}[displayed Energy Readout]
  \label{def:physics:phyx-mini-0582:phyxminiproblems-problemphyxmini0582-displayedenergyreadout}
  \lean{PhyXMiniProblems.ProblemPhyXMini0582.displayedEnergyReadout}
  \uses{def:physics:phyx-mini-0582:phyxminiproblems-problemphyxmini0582-massreadout, def:physics:phyx-mini-0582:phyxminiproblems-problemphyxmini0582-lengthreadout, def:physics:phyx-mini-0582:phyxminiproblems-problemphyxmini0582-planckconstantreadout, def:physics:phyx-mini-0582:phyxminiproblems-problemphyxmini0582-diatomicmolecularrotorsetup, def:physics:phyx-mini-0582:phyxminiproblems-problemphyxmini0582-answerchoice}
  The formula printed beside each answer choice, read in coherent units
\end{definition}
\begin{proof}
  This is the declared model definition of displayed Energy Readout.
\end{proof}
\begin{definition}[recorded Dataset Answer]
  \label{def:physics:phyx-mini-0582:phyxminiproblems-problemphyxmini0582-recordeddatasetanswer}
  \lean{PhyXMiniProblems.ProblemPhyXMini0582.recordedDatasetAnswer}
  \uses{def:physics:phyx-mini-0582:phyxminiproblems-problemphyxmini0582-answerchoice}
  The answer label recorded by the source dataset
\end{definition}
\begin{proof}
  This is the declared model definition of recorded Dataset Answer.
\end{proof}
\begin{definition}[Matches Answer Choice]
  \label{def:physics:phyx-mini-0582:phyxminiproblems-problemphyxmini0582-matchesanswerchoice}
  \lean{PhyXMiniProblems.ProblemPhyXMini0582.MatchesAnswerChoice}
  \uses{def:physics:phyx-mini-0582:phyxminiproblems-problemphyxmini0582-energyreadout, def:physics:phyx-mini-0582:phyxminiproblems-problemphyxmini0582-diatomicmolecularrotorsetup, def:physics:phyx-mini-0582:phyxminiproblems-problemphyxmini0582-answerchoice, def:physics:phyx-mini-0582:phyxminiproblems-problemphyxmini0582-displayedenergyreadout}
  A displayed choice agrees with the modeled energy at level n
\end{definition}
\begin{proof}
  This is the declared model definition of Matches Answer Choice.
\end{proof}
% --- Archon declaration topology end ---
% --- Archon physics formalization source end ---
